\textbf{Lösung zu Übung 4, Aufgabe 1}

\begin{loesung}
$\Omega = \{1 \le k \le 1000\}$. Sei $A_i = \{\ \Omega \mid i \text{ ist Teiler von } 1000\ \}$.
% Hinweis: so im Original (gemeint ist wohl A_i = { k in Omega | i ist Teiler von k }).
Die gesuchte Anzahl ist dann
\begin{align*}
|\,\Omega \setminus A_2 \cup A_3 \cup A_5\,| &= |\Omega| - \alpha_1 + \alpha_2 - \alpha_3
\end{align*}
mit
\begin{align*}
\alpha_1 &= |A_2| + |A_3| + |A_5|\\
\alpha_2 &= |A_2 \cap A_3| + |A_2 \cap A_5| + |A_3 \cap A_5|\\
\alpha_3 &= |A_2 \cap A_3 \cap A_5|
\end{align*}
\begin{align*}
&|A_2| = \frac{1000}{2}, \qquad |A_3| = \left\lfloor \frac{1000}{3} \right\rfloor, \qquad |A_5| = \frac{1000}{5},\\
\Rightarrow\quad &\alpha_1 = 500 + 333 + 200 = 1033\\
&|A_2 \cap A_3| = \left\lfloor \frac{1000}{6} \right\rfloor, \quad |A_2 \cap A_5| = \frac{1000}{10}, \quad |A_3 \cap A_5| = \left\lfloor \frac{1000}{15} \right\rfloor\\
\Rightarrow\quad &\alpha_2 = 166 + 100 + 66 = 332\\
&|A_2 \cap A_3 \cap A_5| = \left\lfloor \frac{1000}{30} \right\rfloor,\\
\Rightarrow\quad &\alpha_3 = 33\\
\Rightarrow\quad &|\,\Omega \setminus A_2 \cup A_3 \cup A_5\,| = |\Omega| - \alpha_1 + \alpha_2 - \alpha_3\\
&= 1000 - 1033 + 332 - 33 = 266
\end{align*}

Schneller findet man die Lösung, wenn man die Produktformel benutzt, die hier aber nur näherungsweise gilt:
\begin{align*}
|\,\Omega \setminus A_2 \cup A_3 \cup A_5\,| &= 1000 \cdot \left(1 - \frac{1}{2}\right)\left(1 - \frac{1}{3}\right)\left(1 - \frac{1}{5}\right)\\
&= 100 \cdot \frac{8}{3}\\
&= 266{,}66
\end{align*}
% Hinweis: "100 * 8/3" so im Original (entspricht 1000 * 4/15).
\end{loesung}
